\section*{Chapter \ref{ch:b1:long-run-model} --- The Long-Run Model}
% Numbers recomputed in code/bachelor-1/19-long-run-model/exercises/solutions.py (tested).

\begin{solution}{exo:b1:long-run-model:prices}
$Y = 10 \times 100^{0.3} \times 400^{0.7} = 2639.02$; $w = 0.7Y/L = 4.618$; $R/P = 0.3Y/K = 7.917$.
\end{solution}

\begin{solution}{exo:b1:long-run-model:read}
Japan's share rose by $7.9$ points from 1970 to 2022; the American share fell
by $6.0$ points from 1950 to 2023.
\end{solution}

\begin{solution}{exo:b1:long-run-model:euler-ex}
$4.618 \times 400 + 7.917 \times 100 = 1847.2 + 791.7 = 2638.9$, equal to $Y = 2639.02$ up to rounding
(exactly, $0.7Y + 0.3Y = Y$).
\end{solution}

\begin{solution}{exo:b1:long-run-model:predict}
\omterm{def:b1:long-run-model:saving}{National saving} falls, so the real interest rate rises and investment falls
by the rise in purchases; consumption does not change, since $Y - T$ is
unchanged.
\end{solution}

\begin{solution}{exo:b1:long-run-model:lf}
$C = 100 + 0.6 \times 1700 = 1120$; \omterm{def:b1:long-run-model:saving}{private saving} $1700 - 1120 = 580$; \omterm{def:b1:long-run-model:saving}{public saving} 0;
\omterm{def:b1:long-run-model:saving}{national saving} 580; $700 - 20r = 580$ gives $r = 6$ and $I = 580$.
\end{solution}

\begin{solution}{exo:b1:long-run-model:taxcut}
$C = 100 + 0.6 \times 1750 = 1150$; $S = 2000 - 1150 - 300 = 550$; $r = (700 - 550)/20 = 7.5$, $I = 550$.
Consumption rises by 30 and investment falls by 30.
\end{solution}

\begin{solution}{exo:b1:long-run-model:immig}
$w \propto (K/L)^{0.3}$ falls by the factor $1.1^{-0.3} = 0.972$ (about $2.8\%$); $R/P \propto (L/K)^{0.7}$ rises
by $1.1^{0.7} = 1.069$. The shares stay at $0.7$ and $0.3$.
\end{solution}

\begin{solution}{exo:b1:long-run-model:trend}
Times 4.3 decades: Germany $-5.59$, France $-4.95$, Japan $-4.34$, United States
$-3.35$ and United Kingdom $+3.23$ points.
\end{solution}

\begin{solution}{exo:b1:long-run-model:brenland}
$393264 - 225014 - 78653 = 89597$, against investment of 88,484. The gap, 1,113, is
net exports $X - M$: Brenland is open, and its saving beyond its investment
is lent abroad (before net income and transfers from abroad).
\end{solution}

\begin{solution}{exo:b1:long-run-model:ces}
With $Y = [aK^{\rho} + (1 - a)L^{\rho}]^{1/\rho}$ and $\sigma = 1/(1 - \rho)$, the ratio of the shares is
$F_K K/(F_L L) = \frac{a}{1 - a}(K/L)^{\rho}$. If $\sigma < 1$, $\rho < 0$, and a rise in $K/L$ lowers capital's
share relative to labour's: labour's share rises. With $\sigma > 1$, $\rho > 0$, and
labour's share falls.
\end{solution}

\begin{solution}{exo:b1:long-run-model:partial}
Write $I = I_0 - i_1r$. \omterm{def:b1:models-and-data:equilibrium}{Equilibrium}: $\bar S - \Delta G + sr = I_0 - i_1r$, so $\Delta r = \Delta G/(s + i_1)$
and $\Delta I = -i_1\Delta G/(s + i_1)$, less than $\Delta G$ in size: the higher rate raises \omterm{def:b1:long-run-model:saving}{private
saving} by $s\Delta r$, which finances part of the purchases.
\end{solution}

\begin{solution}{exo:b1:long-run-model:neutral-ex}
Across countries or decades, a lasting relation between the growth of
money and the level or growth of real output, real wages or real interest
rates that survives controls for other causes --- for example, countries
with faster money growth growing faster or slower for decades for no other
reason.
\end{solution}

\begin{solution}{pb:b1:long-run-model:brenland}
\textbf{1.} $A = 400/(1000^{0.35} \times 4.4^{0.65}) = 13.6087$.

\textbf{2.} $w = 0.65 \times 400/4.4 = 59.09$ thousand dalens.

\textbf{3.} $0.35 \times 400/1000 = 0.14$.

\textbf{4.} $59.09 \times 4.4 + 0.14 \times 1000 = 260 + 140 = 400$.

\textbf{5.} $0.65$ and $0.35$.

\textbf{6.} The wage by $1.1^{-0.35} = 0.9672$ and the rental price by $1.1^{0.65} = 1.0639$.

\textbf{7.} $C = 36 + 0.6 \times 320 = 228$.

\textbf{8.} Private $400 - 80 - 228 = 92$; public $80 - 80 = 0$.

\textbf{9.} 92.

\textbf{10.} $132 - 10r = 92$: $r = 4\%$ and $I = 92$.

\textbf{11.} $228/400 = 0.57$, $92/400 = 0.23$ and $80/400 = 0.2$: close to $0.572$, $0.225$ and $0.2$.

\textbf{12.} $142 - 10r = 92$: $r = 5\%$, with investment unchanged at 92.

\textbf{13.} $\Delta G = 8$: \omterm{def:b1:long-run-model:saving}{public saving} falls to $-8$ and \omterm{def:b1:long-run-model:saving}{national saving} to 84.

\textbf{14.} $132 - 10r = 84$: $r = 4.8\%$.

\textbf{15.} By 8, from 92 to 84.

\textbf{16.} $C = 36 + 0.6 \times 312 = 223.2$; $S = 400 - 223.2 - 88 = 88.8$; $r = (132 - 88.8)/10 = 4.32\%$;
$I = 88.8$.

\textbf{17.} Households pay part of the tax out of saving: consumption falls by
$0.6 \times 8 = 4.8$, so \omterm{def:b1:long-run-model:saving}{national saving} falls by only 3.2.

\textbf{18.} Less investment means a smaller capital stock in future years,
so lower output and, with less capital per worker, lower wages.

\textbf{19.} 440.

\textbf{20.} $w = 0.65 \times 440/4.4 = 65.0$ and $0.35 \times 440/1000 = 0.154$: both up $10\%$.

\textbf{21.} Unchanged at $0.65$ and $0.35$.

\textbf{22.} $C = 36 + 0.6 \times 360 = 252$, $S = 440 - 252 - 80 = 108$, and $r = (132 - 108)/10 = 2.4\%$.

\textbf{23.} A higher $A$ raises the \omterm{def:b1:technology:marginal-product}{marginal product} of capital, so firms
want more capital at every interest rate; investment demand shifts out,
which raises $r$ back up.

\textbf{24.} The defence build-up of $2\%$ of output raises the real interest
rate by $0.8$ points, from $4\%$ to $4.8\%$, and crowds out investment one for one,
by 8 billion dalens.
\end{solution}
